% Options for packages loaded elsewhere
% Options for packages loaded elsewhere
\PassOptionsToPackage{unicode}{hyperref}
\PassOptionsToPackage{hyphens}{url}
\PassOptionsToPackage{dvipsnames,svgnames,x11names}{xcolor}
%
\documentclass[
  11pt,
  letterpaper,
  DIV=11,
  numbers=noendperiod]{scrartcl}
\usepackage{xcolor}
\usepackage{amsmath,amssymb}
\setcounter{secnumdepth}{-\maxdimen} % remove section numbering
\usepackage{iftex}
\ifPDFTeX
  \usepackage[T1]{fontenc}
  \usepackage[utf8]{inputenc}
  \usepackage{textcomp} % provide euro and other symbols
\else % if luatex or xetex
  \usepackage{unicode-math} % this also loads fontspec
  \defaultfontfeatures{Scale=MatchLowercase}
  \defaultfontfeatures[\rmfamily]{Ligatures=TeX,Scale=1}
\fi
\usepackage{lmodern}
\ifPDFTeX\else
  % xetex/luatex font selection
  \setmainfont[]{arial}
\fi
% Use upquote if available, for straight quotes in verbatim environments
\IfFileExists{upquote.sty}{\usepackage{upquote}}{}
\IfFileExists{microtype.sty}{% use microtype if available
  \usepackage[]{microtype}
  \UseMicrotypeSet[protrusion]{basicmath} % disable protrusion for tt fonts
}{}
\makeatletter
\@ifundefined{KOMAClassName}{% if non-KOMA class
  \IfFileExists{parskip.sty}{%
    \usepackage{parskip}
  }{% else
    \setlength{\parindent}{0pt}
    \setlength{\parskip}{6pt plus 2pt minus 1pt}}
}{% if KOMA class
  \KOMAoptions{parskip=half}}
\makeatother
% Make \paragraph and \subparagraph free-standing
\makeatletter
\ifx\paragraph\undefined\else
  \let\oldparagraph\paragraph
  \renewcommand{\paragraph}{
    \@ifstar
      \xxxParagraphStar
      \xxxParagraphNoStar
  }
  \newcommand{\xxxParagraphStar}[1]{\oldparagraph*{#1}\mbox{}}
  \newcommand{\xxxParagraphNoStar}[1]{\oldparagraph{#1}\mbox{}}
\fi
\ifx\subparagraph\undefined\else
  \let\oldsubparagraph\subparagraph
  \renewcommand{\subparagraph}{
    \@ifstar
      \xxxSubParagraphStar
      \xxxSubParagraphNoStar
  }
  \newcommand{\xxxSubParagraphStar}[1]{\oldsubparagraph*{#1}\mbox{}}
  \newcommand{\xxxSubParagraphNoStar}[1]{\oldsubparagraph{#1}\mbox{}}
\fi
\makeatother


\usepackage{longtable,booktabs,array}
\usepackage{calc} % for calculating minipage widths
% Correct order of tables after \paragraph or \subparagraph
\usepackage{etoolbox}
\makeatletter
\patchcmd\longtable{\par}{\if@noskipsec\mbox{}\fi\par}{}{}
\makeatother
% Allow footnotes in longtable head/foot
\IfFileExists{footnotehyper.sty}{\usepackage{footnotehyper}}{\usepackage{footnote}}
\makesavenoteenv{longtable}
\usepackage{graphicx}
\makeatletter
\newsavebox\pandoc@box
\newcommand*\pandocbounded[1]{% scales image to fit in text height/width
  \sbox\pandoc@box{#1}%
  \Gscale@div\@tempa{\textheight}{\dimexpr\ht\pandoc@box+\dp\pandoc@box\relax}%
  \Gscale@div\@tempb{\linewidth}{\wd\pandoc@box}%
  \ifdim\@tempb\p@<\@tempa\p@\let\@tempa\@tempb\fi% select the smaller of both
  \ifdim\@tempa\p@<\p@\scalebox{\@tempa}{\usebox\pandoc@box}%
  \else\usebox{\pandoc@box}%
  \fi%
}
% Set default figure placement to htbp
\def\fps@figure{htbp}
\makeatother





\setlength{\emergencystretch}{3em} % prevent overfull lines

\providecommand{\tightlist}{%
  \setlength{\itemsep}{0pt}\setlength{\parskip}{0pt}}



 


\usepackage{booktabs}
\usepackage{longtable}
\usepackage{array}
\usepackage{multirow}
\usepackage{wrapfig}
\usepackage{float}
\usepackage{colortbl}
\usepackage{pdflscape}
\usepackage{tabu}
\usepackage{threeparttable}
\usepackage{threeparttablex}
\usepackage[normalem]{ulem}
\usepackage{makecell}
\usepackage{xcolor}
\KOMAoption{captions}{tableheading}
\makeatletter
\@ifpackageloaded{caption}{}{\usepackage{caption}}
\AtBeginDocument{%
\ifdefined\contentsname
  \renewcommand*\contentsname{Table of contents}
\else
  \newcommand\contentsname{Table of contents}
\fi
\ifdefined\listfigurename
  \renewcommand*\listfigurename{List of Figures}
\else
  \newcommand\listfigurename{List of Figures}
\fi
\ifdefined\listtablename
  \renewcommand*\listtablename{List of Tables}
\else
  \newcommand\listtablename{List of Tables}
\fi
\ifdefined\figurename
  \renewcommand*\figurename{Figure}
\else
  \newcommand\figurename{Figure}
\fi
\ifdefined\tablename
  \renewcommand*\tablename{Table}
\else
  \newcommand\tablename{Table}
\fi
}
\@ifpackageloaded{float}{}{\usepackage{float}}
\floatstyle{ruled}
\@ifundefined{c@chapter}{\newfloat{codelisting}{h}{lop}}{\newfloat{codelisting}{h}{lop}[chapter]}
\floatname{codelisting}{Listing}
\newcommand*\listoflistings{\listof{codelisting}{List of Listings}}
\makeatother
\makeatletter
\makeatother
\makeatletter
\@ifpackageloaded{caption}{}{\usepackage{caption}}
\@ifpackageloaded{subcaption}{}{\usepackage{subcaption}}
\makeatother
\usepackage{bookmark}
\IfFileExists{xurl.sty}{\usepackage{xurl}}{} % add URL line breaks if available
\urlstyle{same}
\hypersetup{
  colorlinks=true,
  linkcolor={blue},
  filecolor={Maroon},
  citecolor={Blue},
  urlcolor={Blue},
  pdfcreator={LaTeX via pandoc}}


\author{}
\date{}
\begin{document}


\textbf{Minnesotan Crime: Patterns Over Time and Across COVID}

A Capstone Project for the Department of Math \& Statistics

William Hamilton Borer Seabloom

Advised by Dr.Tyler George

Cornell College

Feb x 2026

\subsubsection{Introduction}\label{introduction}

~~~~~ In 1998 J. Mark Norman and Donald E. Arwood {[}1{]} published a
paper titled ``A Social Disorganization Theory of County Crime Rates in
Minnesota'' on how many different sociological measures of a society's
organization and control affect the crime rates in a community. These
factors were measured across Minnesotan counties and included factors
such as education, wages, population density, and police presence.
Recently, during late 2019 and 2020 the COVID-19 pandemic changed many
factors of our society and disrupted many dimensions of social life. As
discussed by Boman and Gallupe in 2020 {[}3{]} one of these many impacts
was on crime rates, particularly in relation to the more serious violent
crimes analyzed in Norman and Arwood's 1998 study {[}1{]}. This raises
an important question: Across Minnesotan counties how have the
relationships between crime and its determinants changed over time,
particularly in response to the COVID-19 pandemic?\\
\strut ~~~~~ The main goal of this research is to build on Norman and
Arwood's 1998 study, by not only redoing the analysis with recent data
but also looking at how these relationships have changed due to the
COVID-19 pandemic. With crime and its prevention being such a large part
of our world, and COVID changing many key parts of our society, it is
important to reevaluate how typical factors believed to determine crime
could have shifted during the COVID-19 pandemic. The state of these
relationships today is also important to evaluate for future
prediction/prevention of crime.\\
\strut ~~~~~ To investigate this question, the study will compare crime
counts from the FBI's Crime Data Explorer with demographic and
socioeconomic indicators drawn from the U.S. Census five-year
county-level estimates. The analysis will include all 87 Minnesota
counties and will operationalize crime rates as the number of Part I
index offenses (homicide, burglary, robbery, aggravated assault, rape,
larceny, motor vehicle theft, and arson) per county, normalized by
population. This definition of crime aligns with Norman and Arwood's
earlier work {[}1{]} while also building on the more serious offenses
that Boman and Gallupe (2020) {[}3{]} identified as trending upward
during the COVID-19 pandemic. Differences in crime rates, averaged
across the five-year periods before and after the onset of the pandemic,
will be examined to assess whether structural changes occurred. These
differences will be evaluated through statistical modeling of both the
pre and post-pandemic county-level data, as well as models estimating
changes over time.

\subsubsection{Literature Review}\label{literature-review}

~~~~~ This project is informed by several key sources: Norman and Arwood
(1998) {[}1{]}, James and Logan (2008) {[}2{]}, and Boman and Gallupe
(2020) {[}3{]}. Together, these works provide a foundation for
understanding the sociological determinants of crime, the limitations of
available data, and the unique disruptions brought about by the COVID-19
pandemic.\\
\strut ~~~~~ Norman and Arwood's 1998 study analyzed the relationship
between crime rates and a variety of sociological measures at the county
level in Minnesota. Their findings demonstrated that structural
indicators such as households missing one or both parents, alcohol
consumption, and migration among other variables are significantly
related to crime. This work provides the methodological and theoretical
foundation for the present project. While their study focused on the
sociological context of the late 20th century, this research builds upon
their approach by incorporating updated data and examining how these
relationships may have shifted due to the COVID-19 pandemic.\\
\strut ~~~~~ James and Logan's 2008 report provides an overview of how
the FBI collects and reports crime statistics, detailing the
methodologies, limitations, and potential sources of bias within Uniform
Crime Reporting (UCR) data. Their report is particularly important
because it highlights challenges such as underreporting, changes in
reporting practices, and differences in local enforcement priorities.
James and Logan state that underreporting in the FBI crime data is a
large problem, with systematic errors produced from the way that the
data is collected. This bias stems from the fact that the only crimes
the FBI report are ones that the police have reported to them, and
police officers are only aware of citizen reported crimes. Changes in
reporting practices also affect this data, in particular the arson and
rape counts. The way that rape was defined in this data changed in 2013
to remove the requirement of it being `forcible' increasing the counts
of this crime around this period as local law enforcement agencies
switched their reporting. Arson counts are also handled differently than
other part 1 index crimes, all other crimes are only recorded as the
most severe offense if any other crimes were committed but arson is
counted on top of any other crime committed. Additionally, arson counts
are only recorded if it can be proven that the fire was willfully set
skewing the counts further.\\
\strut ~~~~~ These insights are crucial for interpreting the results of
this project, as they remind us that observed patterns in county-level
crime rates may partly reflect variation in data quality or reporting
practices rather than actual changes in criminal activity. By
incorporating this perspective, this study can more carefully interpret
both the limitations and strengths of using FBI crime statistics in the
context of this paper's results.\\
\strut ~~~~~ Boman and Gallupe's research examines the effects of the
COVID-19 pandemic on crime trends, with particular attention to how
crime patterns shifted during periods of widespread social disruption.
They found that while many types of crime declined during the pandemic,
homicide and other serious violent crimes rose due to young people being
put in positions where they were particularly close to partners or
family members and further from their friends. This change in social
structures led to less of the smaller crimes such as vandalism but
higher counts of more severe crimes such as homicide and aggravated
assault. Their findings suggest that structural shocks, such as a global
pandemic, can have heterogeneous effects on different categories of
crime. This study directly informs my project, as I will investigate
whether similar shifts occurred in Minnesota counties, particularly in
the prevalence of violent crime relative to overall crime rates as well
as a potential shift in the predictability of crime after the
pandemic.\\
\strut ~~~~~ Taken together, these works establish the intellectual
groundwork for this project. Norman and Arwood demonstrate how
sociological indicators can be systematically linked to crime rates at
the county level. James and Logan's 2008 report provides guidance on the
interpretation of crime statistics, ensuring that my analysis
acknowledges the methodological challenges of working with FBI data.
Finally, Boman and Gallupe highlight the COVID-19 pandemic as a
significant structural event that reshaped crime in the United States,
motivating the inclusion of a temporal comparison between pre and
post-pandemic trends in Minnesota. Building on these sources, my study
aims to contribute by both replicating and extending prior methods to
assess whether longstanding sociological-crime relationships have been
altered by the COVID-19 pandemic.

\subsubsection{Data Description}\label{data-description}

~~~~~ The data being used in this study consists of two main parts, FBI
crime and officer reports as well as U.S Census Bureau county level ACS
5 estimates. A collection of other sources were used to supplement and
condense these main sources as well. All of the data and code for this
project is publicly available in a
\href{https://github.com/lborerseabloom/Econometrics.git}{Github
repository} to document the process of this research. Any graphs,
tables, and results discussed in this paper can be replicated by
following the steps detailed in the repository. ~~~~~ The FBI crime and
officer counts were scraped from their data explorer at an agency level
across all part 1 index crimes. As discussed by James, Nathan, and Logan
Rishard {[}2{]}, this data contains quite a bit of intricacies to be
aware of when using it to draw real world conclusions. Next, using the
data from the NLETS ORI directory {[}4{]} this data was mapped on to and
summarized at the county level.

\pandocbounded{\includegraphics[keepaspectratio]{main_econometrics_files/figure-pdf/homicides_time_series-1.pdf}}

~~~~~ A brief look at the monthly/yearly plot of homicides over time
shows some of the changes in reporting and issues with the data which
are present in all crimes. The most obvious of these are the December
peaks as many police departments report the whole year in December and
the change in reporting practices in 1998 and 2013 where how the FBI
receives and records their data changes James and Logan (2008) {[}2{]}.
However, the grouping of crime data into years removes the month to
month reporting problems and allows for more direct comparison to the
Census estimates.\\
\strut ~~~~~ County level 5 year estimates (ACS5) for the demographic
variables of interest during the years of 2014-2018 and 2019-2023 were
then obtained from the U.S Census Bureau through the tidycensus package.
These estimates have their own problems discussed on the census website
such as sampling and response bias {[}5{]}. The main problems from this
census data are the fact that data on every county is only available in
5 year average estimates and that these numbers are estimates based on a
sample of the population. However, the latter is not a problem in our
case as we want to examine structural changes resulting from the
COVID-19 pandemic and an average across 5 years will be able to
represent that. The former nonetheless, will have to be addressed in the
final interpretation of the results as a potential source of bias.\\
\strut ~~~~~ With this data the final dependent variable of interest,
crime rate, is calculated by summing a 5 year average of all part 1
index crimes and dividing it by the total population estimate from the
census. The demographic variables used in Norman and Arwood's 1998 study
{[}1{]} were then calculated using a variety of Census estimates. There
was no measure of alcohol taxes by county as used in their original
study, {[}1{]} but the Minnesota department of health published a report
on the prevalence of binge drinking rates by county {[}6{]} which can be
used as a similar indicator for alcohol use. An additional factor found
to be significant in the original study, {[}1{]} migration estimates,
were unavailable due to the lack of federal funding for some census
estimates.\\
\strut ~~~~~ The final data being used for modeling contains two rows
for each of the 87 Minnesota counties representing the pre and post
COVID periods with 5 year average estimates for 2014-2018 and 2019-2023
respectively. The variables used in Norman and Arwood's 1998 study
{[}1{]} will be tested in their effects on crime rates before and after
to address the question of: how have the relationships between crime and
its determinants changed over time, particularly in response to the
COVID-19 pandemic?

\subsubsection{Research Findings}\label{research-findings}

~~~~~ Using this data a few graphical representations of crime rates can
help to better understand the nature of the data before modeling.

\pandocbounded{\includegraphics[keepaspectratio]{main_econometrics_files/figure-pdf/homicides_time_series2-1.pdf}}

~~~~~ Using ACS 1 year estimates of population, a yearly plot of crime
rates and population reveals a consistent downward trend in crime rates
per capita, dropping with a steeper slope than the increase in
population. This pattern suggests that crime rates have been declining
naturally over time. However, consistent with the predictions in the
study by Boman, J.H., Gallupe, O, {[}3{]} a notable spike in crime is
observed between 2019 and 2020, coinciding with the onset of the
COVID-19 pandemic. Following this deviation, crime rates decline again
by 2023, resuming the pre-pandemic trend. This indicates that the
pandemic period represents a temporary disturbance rather than a lasting
structural shift in Minnesotan crime trends.

\pandocbounded{\includegraphics[keepaspectratio]{main_econometrics_files/figure-pdf/homicides_time_series3-1.pdf}}

~~~~~ Looking at the general distribution of crime rates we see a fairly
normal distribution for the majority of the counties with a right tail
holding the largest counties. In this plot we can clearly see the
reduction in crime rates in Minnesotan counties as well as the
correlation between population centers and crime.

\pandocbounded{\includegraphics[keepaspectratio]{main_econometrics_files/figure-pdf/crime_and_pop_mapping-1.pdf}}

~~~~~ A spatial analysis of the 2023 post-COVID five-year averages for
crime rates and population reveals distinct geographic patterns. Higher
crime rates are concentrated within major urban centers, with a
noticeable decline in surrounding suburban areas. Outside of the Twin
Cities region, other population hubs also exhibit above-average crime
rates. One clear outlier is Beltrami County, home to the Red Lake
Reservation, which can be seen as a clear yellow island in the upper
central portion of the crime rate map. It should be noted that the
``percent non-white'' variable partly captures demographic patterns
associated with such reservations.

\begin{longtable}[]{@{}lrrrrr@{}}
\caption{Signifigance of changes in Cesus variables}\tabularnewline
\toprule\noalign{}
variable & mean\_pre & mean\_post & sd\_pre & sd\_post & p.value \\
\midrule\noalign{}
\endfirsthead
\toprule\noalign{}
variable & mean\_pre & mean\_post & sd\_pre & sd\_post & p.value \\
\midrule\noalign{}
\endhead
\bottomrule\noalign{}
\endlastfoot
median\_household\_income\_acs5 & 5.956394e+04 & 7.575739e+04 &
1.095935e+04 & 1.304152e+04 & 0.0000000 \\
pct\_lessthan\_5kincome & 4.777142e+00 & 3.879403e+00 & 1.362166e+00 &
1.143885e+00 & 0.0000052 \\
pct\_non\_white & 8.670894e+00 & 1.271564e+01 & 7.473397e+00 &
8.355481e+00 & 0.0009449 \\
pct\_highschool\_or\_greater & 9.109053e+01 & 9.227239e+01 &
2.706797e+00 & 2.430826e+00 & 0.0028271 \\
pct\_children\_missing\_parents & 3.337682e+01 & 3.244493e+01 &
1.891822e+00 & 2.498210e+00 & 0.0061994 \\
crime\_rate & 1.452060e-02 & 1.126310e-02 & 8.316000e-03 & 7.683000e-03
& 0.0079983 \\
pct\_young\_males & 3.194073e+00 & 3.365251e+00 & 5.433701e-01 &
5.733391e-01 & 0.0448063 \\
officers & 2.751613e+00 & 2.702441e+00 & 9.525891e-01 & 7.305777e-01 &
0.7029255 \\
persons\_per\_household & 2.462972e+00 & 2.466778e+00 & 2.006622e-01 &
1.643630e-01 & 0.8912944 \\
persons\_per\_m2 & 4.940000e-05 & 5.090000e-05 & 1.608000e-04 &
1.634000e-04 & 0.9490400 \\
Binge.drinking.prevalence & 2.319651e+01 & 2.319651e+01 & 1.299995e+00 &
1.299995e+00 & 1.0000000 \\
\end{longtable}

~~~~~ An examination of the predictive variables used in {[}1{]}
indicates no statistically significant changes at the 5\% level between
the pre and post-COVID periods for officers per 1,000 residents, persons
per household, and persons per square meter. The binge drinking
prevalence variable here, is a point-estimate and therefore remains
constant over time, as expected.

\begin{longtable}[]{@{}lrr@{}}
\caption{Signifigance of correlation with crime rates}\tabularnewline
\toprule\noalign{}
variable & estimate & p.value \\
\midrule\noalign{}
\endfirsthead
\toprule\noalign{}
variable & estimate & p.value \\
\midrule\noalign{}
\endhead
\bottomrule\noalign{}
\endlastfoot
persons\_per\_m2 & 0.4696967 & 0.0000000 \\
pct\_non\_white & 0.4532249 & 0.0000000 \\
Binge.drinking.prevalence & -0.3263230 & 0.0000125 \\
persons\_per\_household & 0.2595443 & 0.0005434 \\
pct\_lessthan\_5kincome & 0.1924299 & 0.0109639 \\
pct\_young\_males & 0.1771871 & 0.0193380 \\
pct\_highschool\_or\_greater & -0.1715511 & 0.0236104 \\
officers & 0.1292337 & 0.0892144 \\
pct\_children\_missing\_parents & 0.0856260 & 0.2612643 \\
median\_household\_income\_acs5 & -0.0228655 & 0.7645723 \\
\end{longtable}

~~~~~ Further correlation analyses show that while several predictors
exhibit meaningful and statistically significant relationships with
crime rates, there is little evidence of substantial change during the
COVID-19 pandemic period. Interestingly, the variables with the least
variation pre and post COVID are those most strongly correlated with
crime rates at the county level.\\
\strut ~~~~~ A full linear regression model confirms these patterns.
Several predictors were excluded due to statistical insignificance (p
\textgreater{} 0.5), including percent of children missing one or both
parents, msp\_main\_counties: and indicator for if the county is Ramsey
or Hennepin the two main counties for Saint Paul and Minneapolis,
Binge.drinking.prevalence, pct\_young\_males: the percent of the male
population between 15 and 25, and persons per household. The percent
less than \$5000 income variable was also removed due to its high
negative correlation (r = -0.7) with median income and its theoretical
incorporation of much of the same information.

\begin{table}

\caption{\label{tab:acs5_modeling}Regression Models with Classical and Robust Standard Errors (HC3)}
\centering
\begin{tabular}[t]{l|l|r|r|r|r}
\hline
Model & term & estimate & std.error\_normal & std.error\_robust & p.value\\
\hline
\multicolumn{6}{l}{\textbf{Pre-COVID}}\\
\hline
\hspace{1em}Pre-COVID & (Intercept) & 0.0274983 & 0.026 & 0.023 & 0.226\\
\hline
\hspace{1em}Pre-COVID & log\_officers & -0.0014497 & 0.002 & 0.002 & 0.523\\
\hline
\hspace{1em}Pre-COVID & median\_household\_income\_acs5 & -0.0000005 & 0.000 & 0.000 & 0.000\\
\hline
\hspace{1em}Pre-COVID & pct\_highschool\_or\_greater & 0.0006223 & 0.000 & 0.000 & 0.002\\
\hline
\hspace{1em}Pre-COVID & log\_pct\_non\_white & 0.0048612 & 0.001 & 0.001 & 0.000\\
\hline
\hspace{1em}Pre-COVID & log\_persons\_per\_m2 & 0.0045094 & 0.001 & 0.001 & 0.000\\
\hline
\hspace{1em}Pre-COVID & reservations & 0.0042729 & 0.002 & 0.002 & 0.027\\
\hline
\hspace{1em}Pre-COVID & top\_cities & 0.0036786 & 0.003 & 0.004 & 0.364\\
\hline
\multicolumn{6}{l}{\textbf{Post-COVID}}\\
\hline
\hspace{1em}Post-COVID & (Intercept) & 0.0355120 & 0.029 & 0.025 & 0.160\\
\hline
\hspace{1em}Post-COVID & log\_officers & 0.0031652 & 0.002 & 0.002 & 0.133\\
\hline
\hspace{1em}Post-COVID & median\_household\_income\_acs5 & -0.0000003 & 0.000 & 0.000 & 0.000\\
\hline
\hspace{1em}Post-COVID & pct\_highschool\_or\_greater & 0.0004012 & 0.000 & 0.000 & 0.102\\
\hline
\hspace{1em}Post-COVID & log\_pct\_non\_white & 0.0027439 & 0.001 & 0.002 & 0.135\\
\hline
\hspace{1em}Post-COVID & log\_persons\_per\_m2 & 0.0045913 & 0.001 & 0.001 & 0.000\\
\hline
\hspace{1em}Post-COVID & reservations & 0.0040764 & 0.001 & 0.001 & 0.003\\
\hline
\hspace{1em}Post-COVID & top\_cities & 0.0079771 & 0.003 & 0.003 & 0.016\\
\hline
\multicolumn{6}{l}{\textbf{Full Model}}\\
\hline
\hspace{1em}Full Model & (Intercept) & 0.0326651 & 0.019 & 0.015 & 0.028\\
\hline
\hspace{1em}Full Model & log\_officers & 0.0006537 & 0.002 & 0.002 & 0.704\\
\hline
\hspace{1em}Full Model & median\_household\_income\_acs5 & -0.0000003 & 0.000 & 0.000 & 0.000\\
\hline
\hspace{1em}Full Model & pct\_highschool\_or\_greater & 0.0004661 & 0.000 & 0.000 & 0.001\\
\hline
\hspace{1em}Full Model & log\_pct\_non\_white & 0.0039591 & 0.001 & 0.001 & 0.000\\
\hline
\hspace{1em}Full Model & log\_persons\_per\_m2 & 0.0044178 & 0.001 & 0.001 & 0.000\\
\hline
\hspace{1em}Full Model & reservations & 0.0043436 & 0.001 & 0.001 & 0.000\\
\hline
\hspace{1em}Full Model & top\_cities & 0.0060426 & 0.002 & 0.002 & 0.009\\
\hline
\hspace{1em}Full Model & post\_covid & -0.0002185 & 0.001 & 0.001 & 0.842\\
\hline
\end{tabular}
\end{table}

~~~~~ The resulting model reveals multiple noteworthy relationships
between county-level crime and the sociological factors outlined in
Norman, J. Mark and Arwood, Donald E's study {[}1{]}. This model also
incorporates the new features: reservations: an indicator for if a
county has a reservation in it, top\_cities: an indicator for Ramsey,
Hennipen, Olmsted, St.~Louis, and Stearns counties which are counties
containing the 5 largest cities in Minnesota, and a pre and post COVID
indicator in the full model. The reservations and cities can be seen in
the previous crime rates map.

\begin{table}

\caption{\label{tab:acs5_modeling_performance}Model Performance Summary (R², Adjusted R², RMSE)}
\centering
\begin{tabular}[t]{l|r|r|r}
\hline
Model & R2 & Adj\_R2 & RMSE\\
\hline
Pre-COVID & 0.639 & 0.607 & 0.005\\
\hline
Post-COVID & 0.657 & 0.626 & 0.004\\
\hline
Full Model & 0.641 & 0.624 & 0.005\\
\hline
\end{tabular}
\end{table}

~~~~~ These models do a decent job at predicting crime at the county
level with only small variation between each model's predictive power.
Heteroskedastic robust errors are used to compensate for some of the
variance issues, particularly with medium level crime rates seen in this
residuals vs fits plot. In order to fix this a log transformation on
crime rate was tried to no effect.

\pandocbounded{\includegraphics[keepaspectratio]{main_econometrics_files/figure-pdf/acs5_modeling_resids_v_fits-1.pdf}}

~~~~~ The first notable finding in the current model concerns median
household income, which shows a strong and consistent negative
relationship with crime rates across all specifications. Counties with
higher median income experience lower levels of crime, and this effect
remains stable both before and after the onset of COVID-19. This result
supports the idea that people with higher income are less likely to
commit any of the part 1 index crimes. In Norman, J. Mark and Arwood,
Donald E's study, the related variable ``percent below poverty'' was
positively associated with crime, although only marginally significant.
Because income and poverty are inverse measures of economic well-being,
both studies findings are consistent. They indicate that economic
deprivation remains one of the most powerful predictors of crime. The
stronger and more consistent effect observed in the present analysis
suggests that the relationship between income and crime may have become
more pronounced over time.\\
\strut ~~~~~ The racial composition (heterogeneity) of counties also
shows a strong relationship with crime in the current model. The log of
the percent of non-white residents is positively and significantly
associated with crime rates in the total and pre-COVID models but
interestingly, not in the post COVID model. This suggests that counties
with higher proportions of non-white populations tended to have higher
recorded crime before COVID but after COVID this effect has dropped.
This change in relationship likely reflects broader geographic and
demographic patterns rather than a direct causal effect. In particular,
areas with higher non-white populations like urban centers and
reservations, also tend to be areas with more concentrated economic
disadvantage and greater population density. For reservations there is a
consistent large and significant effect without much change due to
covid. What we observe in the indicator for large cities is also a
positive estimate with a large increase in significance and effect after
COVID, which could be representative of a change in how crimes are
reported or committed. Norman, J. Mark and Arwood, Donald E's study used
two demographic diversity indicators, heterogeneity and migration rates.
Those variables were only found to be weakly significant in the earlier
study, {[}1{]} suggesting that demographic clustering has become a more
prominent predictor of crime in recent years.\\
\strut ~~~~~ Educational attainment also demonstrates an interesting
shift. In the current model, the percentage of the population with a
high school education or greater is positively and significantly related
to crime. This finding appears counter intuitive since education is
typically viewed as a protective factor. One likely explanation is that
higher education levels are associated with more urban and densely
populated counties, which tend to have greater recorded crime rates due
to both higher population density and higher reporting rates. In
contrast, Norman, J. Mark and Arwood, Donald E's study found that
education variables were statistically insignificant {[}1{]} which is
what is found in the second, post-COVID model. This difference may
indicate that, over time, education has become more of a proxy for
urbanization rather than a direct measure of social stability and that
the more significant large cities variable, like with heterogeneity,
explains this better post-COVID.\\
\strut ~~~~~ Population density is another variable that emerges as an
important and consistent predictor of crime. In all versions of the
current model, the log of persons per square meter is positively and
strongly significant. This indicates that denser counties experience
higher crime rates. The stability of this relationship before and after
COVID-19 suggests that density-related factors, such as increased
opportunity and social strain remain key drivers of crime. In Norman, J.
Mark and Arwood, Donald E's study, population density was represented by
population per mile, which was only marginally significant {[}1{]}. The
present study therefore provides stronger evidence that population
concentration is a reliable and robust determinant of crime.\\
\strut ~~~~~ Law enforcement presence shows a little more nuanced
relationship. Before COVID, the number of officers per capita had no
significant relationship with crime. After COVID, however, the
coefficient becomes positive and slightly more significant, suggesting
that counties with more officers per capita might have had slightly
higher recorded crime rates. This may not indicate an increase in crime
itself but rather greater enforcement activity or reporting intensity in
areas with more officers. Norman, J. Mark and Arwood, Donald E's study
included a similar variable of the police to citizen ratio in a county,
also finding it insignificant in its effects {[}1{]}. The lack of
significance found from 1998 to present and across the COVID-19 pandemic
is the opposite of what one would expect from an indicator of police
presence which theoretically should be more highly correlated with
crime.\\
\strut ~~~~~ Looking at the full model, the post-COVID indicator
variable is negative, but not statistically significant. This suggests
that, once other socioeconomic and demographic factors are accounted
for, overall crime levels did not change significantly in the years
following the pandemic. This finding is reinforced by chow and cusum
tests that found no evidence of structural change over time as seen
below.

\begin{verbatim}

    M-fluctuation test

data:  total_model
f(efp) = 1.3747, p-value = 0.3434
\end{verbatim}

\begin{verbatim}

    OLS-based CUSUM test

data:  cumsum_process
S0 = 0.57153, p-value = 0.8996
\end{verbatim}

~~~~~ Fitting a model using the differences between pre and post covid
to see if any of the changes in explanatory variables might have had an
effect on the change in crime rate, the only variable found to have a
statistically significant effect on the change in crime rates is the
logarithm of the percentage of the population identified as white. This
finding further supports the conclusion that most sociological
predictors of crime remained stable throughout the COVID-19 pandemic
period but the heterogeneity of the population had significant
correlation with the difference. The significance of the percent
non-white variable could again be representative of the areas which were
most affected by the pandemic, such as urban centers and reservations.

\begin{table}

\caption{\label{tab:differenced_data}Regression Model for Differences in Crime Rate}
\centering
\begin{tabular}[t]{l|l|r|r|r}
\hline
Model & term & estimate & std.error & p.value\\
\hline
\multicolumn{5}{l}{\textbf{Difference Model}}\\
\hline
\hspace{1em}Differenced Model & (Intercept) & -0.0062571 & 0.002 & 0.003\\
\hline
\hspace{1em}Differenced Model & log\_officers\_diff & -0.0022757 & 0.003 & 0.475\\
\hline
\hspace{1em}Differenced Model & median\_household\_income\_acs5\_diff & 0.0000001 & 0.000 & 0.480\\
\hline
\hspace{1em}Differenced Model & pct\_lessthan\_5kincome\_diff & -0.0001279 & 0.000 & 0.696\\
\hline
\hspace{1em}Differenced Model & pct\_highschool\_or\_greater\_diff & -0.0000861 & 0.000 & 0.789\\
\hline
\hspace{1em}Differenced Model & log\_pct\_non\_white\_diff & 0.0040791 & 0.001 & 0.005\\
\hline
\hspace{1em}Differenced Model & log\_persons\_per\_m2\_diff & -0.0134346 & 0.013 & 0.321\\
\hline
\end{tabular}
\end{table}

~~~~~ The findings from this study show that the main factors shaping
crime in Minnesota have stayed largely consistent over time, even
through the social disruptions due to the COVID-19 pandemic. While there
were temporary changes in overall crime levels during the pandemic
years, the underlying relationships between crime and its economic,
demographic, and geographic predictors remain stable.

\begin{table}

\caption{\label{tab:lasso_regression}LASSO Regression Summary}
\centering
\begin{tabular}[t]{l|l}
\hline
  & s0\\
\hline
(Intercept) & 0.187371901290471\\
\hline
log\_persons\_per\_m2 & 0.00457759112018756\\
\hline
log\_pct\_non\_white & 0.00289217491746882\\
\hline
reservations & 0.00228809706102866\\
\hline
log\_officers & 0.00129977535651159\\
\hline
top\_cities & 0.000970885037473073\\
\hline
pct\_lessthan\_5kincome & 0.000251117400679103\\
\hline
pct\_highschool\_or\_greater & 0.000123414924227411\\
\hline
pct\_non\_white & 2.60001279935491e-05\\
\hline
male\_20\_24\_acs5 & 1.28809151019549e-07\\
\hline
area\_m2 & 1.02188118677755e-12\\
\hline
median\_household\_income\_acs5 & -2.46390653819108e-07\\
\hline
year & -6.41161302242918e-05\\
\hline
post\_covid & -8.61648422877703e-05\\
\hline
officers & .\\
\hline
male\_15\_19\_acs5 & .\\
\hline
total\_population\_acs5 & .\\
\hline
white\_population\_acs5 & .\\
\hline
children\_total\_acs5 & .\\
\hline
children\_w\_mother\_only\_acs5 & .\\
\hline
children\_w\_father\_only\_acs5 & .\\
\hline
children\_no\_parents\_acs5 & .\\
\hline
edu\_total\_25plus\_acs5 & .\\
\hline
edu\_hs\_grad\_acs5 & .\\
\hline
edu\_some\_college\_lt1yr\_acs5 & .\\
\hline
edu\_some\_college\_gt1yr\_acs5 & .\\
\hline
edu\_assoc\_degree\_acs5 & .\\
\hline
edu\_bachelors\_acs5 & .\\
\hline
edu\_masters\_acs5 & .\\
\hline
edu\_professional\_degree\_acs5 & .\\
\hline
edu\_doctorate\_acs5 & .\\
\hline
households\_total\_acs5 & .\\
\hline
households\_income\_lt5k\_acs5 & .\\
\hline
pct\_children\_missing\_parents & .\\
\hline
pct\_young\_males & .\\
\hline
persons\_per\_household & .\\
\hline
persons\_per\_m2 & .\\
\hline
\end{tabular}
\end{table}

~~~~~ Running a LASSO regression model we can further see the relative
importance of each variable in our model. Included are every relevant
variable used to create the data, as can be seen the variables selected
as reasonably important match up well with the ones manually selected
through relevance and significance in the previous modeling steps.

\begin{table}

\caption{\label{tab:lasso_regression_metrics}LASSO Model Performance Metrics}
\centering
\begin{tabular}[t]{l|r}
\hline
Metric & Value\\
\hline
Best Lambda & 0.0001612\\
\hline
R-squared & 0.6885289\\
\hline
Adjusted R-squared & 0.6632218\\
\hline
RMSE & 0.0045343\\
\hline
\end{tabular}
\end{table}

~~~~~ Additionally we see no meaningful improvement in across any
metrics in using LASSO for our model and the very small value selected
for the best lambda shows that the OLS approach we were using previously
is going to be a good choice.

\subsubsection{Conclusions}\label{conclusions}

~~~~~ Across all models, median household income continues to be one of
the strongest predictors of crime. Counties with higher income levels
consistently show lower crime rates, supporting the long-standing idea
that economic well-being reduces criminal activity. This relationship is
even stronger today than it was in Norman and Arwood's 1998 study,
{[}1{]} suggesting that economic inequality has become a more powerful
influence on community safety.\\
\strut ~~~~~ Population density also remains a key factor. More densely
populated counties tend to have higher crime rates, both before and
after the COVID-19 pandemic. This likely reflects the social and
environmental pressures that come with larger populations such as higher
opportunity for crime and more social tension. The consistency of this
result suggests that urbanization continues to play a central role in
shaping where and how crime occurs.\\
Nevertheless, some relationships have shifted in interesting ways.
Education, for example, is only positively related to crime pre-COVID,
which at first seems counterintuitive. This is probably because counties
with higher levels of education also tend to be more urban, and urban
areas naturally have higher reported crime rates due to density and
better reporting systems. In this case, education may be capturing the
effects of urban living rather than reflecting a direct link between
schooling and crime which could also be the reason for the drop in
significance post-COVID.\\
\strut ~~~~~ The racial composition of counties also changed in its
relationship with crime before and after COVID. Pre-COVID, counties with
higher percentages of non-white residents tended to have higher crime
rates. After COVID, this relationship weakens and is no longer
statistically significant. This could mean that the link between race
and recorded crime lessened during the COVID-19 pandemic, possibly due
to changes in mobility, enforcement, or reporting practices. It might
also reflect the different ways communities were affected by the
pandemic, especially in urban centers and areas with reservations.\\
\strut ~~~~~ Law enforcement presence sees some slight pre and post
pandemic changes as well. Before COVID and after, the number of officers
per capita was not significantly related to crime. After the pandemic,
however, the relationship became slightly positive but still not
significant. If this increase in effect is real, this would suggest that
counties with more officers also saw higher recorded crime rates.
Theoretically, one would expect this to have a higher effect due to
better reporting, data collection, and more crime where police resources
are concentrated.\\
\strut ~~~~~ There are also some limitations that affect how these
results should be interpreted. The FBI's crime data and the U.S. Census
Bureau's ACS estimates both have known problems. FBI data are affected
by under reporting, differences in how agencies define and record
crimes, and changes in reporting systems that happened around the same
time as COVID. The ACS data provide five-year averages rather than
yearly snapshots, which smooths out short-term changes and makes it
harder to isolate the immediate impact of the COVID-19 pandemic. The
models also show signs of uneven variance, even after using robust
standard errors, suggesting that certain counties have unique local
factors not captured by the available data.\\
Future work could expand this analysis in several ways. From a
geo-spatial perspective, models could help capture how crime in one
county influences its neighbors, especially around the Twin Cities and
other urban areas. Using smaller geographic units like census tracts
could reveal more detailed local patterns. Over time, more granular
yearly data could allow for time-series or panel analyses to see how
relationships change year to year, rather than only comparing pre and
post-COVID averages.\\
\strut ~~~~~ It would also be valuable to bring in new types of data.
Measures of housing instability, unemployment fluctuations, or community
engagement could help explain why some counties deviate from expected
trends. Particularly, obtaining an estimate of net migration which was
found highly significant in Norman, J. Mark and Arwood, Donald E's study
but not currently available from the U.S Census. Including these newer
variables could help capture shifts in social behavior that the
variables used in this study miss.\\
\strut ~~~~~ In summary, this study finds that while the COVID-19
pandemic caused a short-term disruption in crime, the broader structural
relationships between crime and its main determinants: income,
population density, and demographic composition have held steady.
Economic inequality remains one of the most important predictors of
crime, while urbanization continues to shape where it happens. Education
and race have become more complex indicators that now seem to reflect
the changing geography and social makeup of Minnesota's counties rather
than direct causes of crime. Although these results suggest that some
social forces driving crime are deeply rooted and resilient, some forces
such as heterogeneity of a community can be shifted in the face of major
global events like the COVID-19 pandemic.

\textbf{Citations}

{[}1{]} Norman, J. Mark \& Arwood, Donald E. (1998). \emph{A Social
Disorganization Theory of County Crime Rates in Minnesota}. \emph{Great
Plains Sociologist}, 11(2), Article 2.\\
\url{https://openprairie.sdstate.edu/greatplainssociologist/vol11/iss2/2}

{[}2{]} James, Nathan \& Logan, Rishard. (2008). \emph{How Crime in the
United States Is Measured}. Congressional Research Service, Library of
Congress.\\
\url{https://apps.dtic.mil/sti/html/tr/ADA476144/}

{[}3{]} Boman, J. H., \& Gallupe, O. (2020). \emph{Has COVID-19 Changed
Crime? Crime Rates in the United States during the Pandemic}.
\emph{American Journal of Criminal Justice}, 45, 537--545.\\
\url{https://doi.org/10.1007/s12103-020-09551-3}

{[}4{]} National Law Enforcement Telecommunications System (NLETS).
(1981). \emph{ORI (Originating Agency Identifier) Directory}. U.S.
Department of Justice. Grant No.~75-SS-99-6018.

{[}5{]} U.S. Census Bureau. (2025, August 19). \emph{Comparing ACS
Data}.\\
\url{https://www.census.gov/programs-surveys/acs/guidance/comparing-acs-data.html}

{[}6{]} Minnesota Department of Health. (n.d.). \emph{Alcohol Costs by
County} {[}PDF{]}.\\
\url{https://www.health.state.mn.us/communities/alcohol/documents/alcoholcostscounty.pdf}




\end{document}
